% Bagian: sets-functions-relations
% Bab: size-of-sets-complete

\documentclass[../../../include/open-logic-chapter]{subfiles}

\begin{document}

\olchapter{sfr}{siz}{Ukuran Himpunan}

\begin{editorial}
Bab ini membahas enumerasi, keterhitungan, dan ketakterhitungan.
Beberapa bagian tersedia dalam dua versi: versi yang lebih elementer,
yang memandang enumerasi sebagai daftar atau surjeksi dari $\PosInt$; dan
versi yang lebih abstrak, yang mendefinisikan enumerasi sebagai bijeksi
dengan $\Nat$.
\end{editorial}

\olimport{introduction}

\olimport{enumerability}

\olimport{zig-zag}

\olimport{pairing}

\olimport{pairing-alt}

\olimport{non-enumerability}

\olimport{reduction}

\olimport{equinumerous-sets}

\olimport{comparing-size}

\olimport{schroder-bernstein}

\begin{editorial}
\olref[sfr][siz][enm-alt]{sec},
\olref[sfr][siz][nen-alt]{sec}, dan \olref[sfr][siz][red-alt]{sec}
berikut merupakan versi alternatif dari \olref[sfr][siz][enm]{sec},
\olref[sfr][siz][nen]{sec}, dan \olref[sfr][siz][red]{sec} yang ditulis
oleh Tim Button untuk digunakan dalam teks Open Set Theory karyanya.
Versi-versi tersebut sedikit lebih lanjut dan menggunakan definisi
enumerabilitas yang berbeda, yang lebih cocok dalam konteks teori himpunan
(yakni, bijeksi dengan $\Nat$ atau suatu segmen awal, alih-alih dapat
dicantumkan dalam daftar atau menjadi daerah hasil suatu fungsi surjektif
dari $\PosInt$).
\end{editorial}

\olimport{enumerability-alt}
\olimport{non-enumerability-alt}
\olimport{reduction-alt}

\OLEndChapterHook

\end{document}
